\section{Reading the clock off prices}
\label{sec:clkinverse}

Everything so far consumed a calendar: someone typed in ``four extra
days on August~26,'' and the clock did the rest.  For scheduled
earnings that is the right workflow --- the dates are published, and
a desk should use them.  But the sizes $N_e$ are nobody's published
number, and the board's own prices carry an opinion about them: we
\emph{saw} that opinion in \cref{fig:clkboard}(b), where one interval
accrued \MacClkBoardHotOverLull$\times$ the variance per day of its
neighbour.  This section turns that observation into an estimator ---
read the calendar \emph{off} the prices --- and, just as importantly,
measures what such an estimator can and cannot deliver.

Set up the objects.  The data is the quoted ladder: expiries
$t_1<\cdots<t_n$ with at-the-money total variances $w_1,\dots,w_n$
(the same per-interval increments $\Delta w_i=w_i-w_{i-1}$,
$\Delta t_i=t_i-t_{i-1}$ as before, with $w_0=t_0=0$).  The unknowns
are candidate event sizes: one candidate per interval, $N_i\ge0$
extra days somewhere inside interval $i$ --- where exactly inside
turns out not to matter, and a proposition below says why.  Given a
candidate vector $N=(N_1,\dots,N_n)$, each interval's width on the
variance clock is
\[
\Delta\tau_i=\Delta t_i+\frac{N_i}{365},
\]
and the chapter's central diagnostic gets its symbol: the
\emph{forward variance} of interval $i$,
\[
\mathcal{F}_i=\frac{\Delta w_i}{\Delta\tau_i}
\]
--- variance accrued per variance-clock year between two consecutive
expiries.  (Calligraphic $\mathcal{F}$; the roman $F$ remains the
forward price of Chapter~6.)  Two structural facts, each one line of
algebra, shape everything.

\emph{First: the data half of $\mathcal{F}_i$ is untouchable.}  The
increments $\Delta w_i$ are fixed by prices
(\cref{prop:clkgauge}(i)), and the calendar widths $\Delta t_i$ by
the listing.  The candidate $N$ enters only through the denominators
--- an estimator of the clock adjusts how the accrual is
\emph{divided}, never how much accrual there \emph{is}.

\emph{Second: an event can only pull its own interval down.}
Interval $j$'s width $\Delta\tau_j$ contains only interval $j$'s own
candidate --- an event in interval $i$ shifts the running clock
$\tau$ of every later expiry equally at both endpoints, so every
other width is unchanged.  And $N_i$ can only \emph{increase}
$\Delta\tau_i$, hence only \emph{decrease} $\mathcal{F}_i$.  The
estimator can lower a spike toward its neighbours; it cannot raise a
valley.  Its entire vocabulary is ``this interval's days were worth
more than ordinary days.''

Now the estimator itself.  On the market's own clock, accrual should
be steady (\cref{sec:clkwalk}); so ask for the candidate vector that
makes the $\mathcal{F}_i$ ladder \emph{as flat as possible} --- while
insisting, because many candidate vectors flatten almost equally
well, that events be few and small.  Formally, minimize over
$N\ge0$:
\begin{equation}
J(N)=\sum_{i<n}\big(\mathcal{F}_{i+1}-\mathcal{F}_i\big)^2
+\lambda_{\rm mono}\sum_{i<n}
\big(\!\min(\mathcal{F}_{i+1}-\mathcal{F}_i,\,0)\big)^2
+\lambda_{\rm sparse}\sum_i\frac{N_i}{365}
+\lambda_{\rm ridge}\sum_i\Big(\frac{N_i}{365}\Big)^{2}.
\label{eq:clkobj}
\end{equation}
Read the four terms one at a time.  The first is the flatness ideal:
squared steps of the forward-variance ladder.  The second charges
\emph{decreases} again, on top.  Why single out decreases?  Because a
genuine pre-event board leaves a specific signature: the event
interval spikes above both neighbours, so the ladder rises into it
and falls out of it.  Lowering that one interval repairs both edges
at once.  An ordinary \emph{upward-sloping} term structure, by
contrast, has no falling edges at all --- markets show that shape
constantly, no event caused it, and the extra charge makes sure the
solver is not tempted to ``repair'' it.  The third term is the cost
of speech: a charge on each installed day, linear in $N_i$.  Linear,
not squared, and the difference matters.  Chapter~2's calibration
carried a \emph{ridge} penalty --- a squared-size charge --- and a
squared charge shrinks everything a little while zeroing nothing;
a linear charge (with $N_i\ge0$ it is exactly the absolute-value
penalty statisticians call $\ell_1$, or \emph{sparsity}) prefers
exact zeros, paying for an event only where the flatness gain clears
a threshold.  The last term is a tiny squared charge that keeps the
optimization well conditioned; it never decides anything.  The minimization is a small bounded solve ---
$n$ variables, reference weights and the routine in
appendix~\ref{app:clkprotocol} --- and solved events below half a day
are dropped as noise.

Before trusting any output, be precise about what the data can
identify at all.

\begin{proposition}[Dates inside an interval are not identified]
\label{prop:clkident}
Every quantity the pipeline consumes --- fits, readings,
interpolation, the objective \eqref{eq:clkobj} --- depends on the
event calendar only through the clock values $\tau(t_1),\dots,
\tau(t_n)$ at the quoted expiries, hence only through the per-interval
totals $N_i=\sum_{t_e\in(t_{i-1},t_i]}N_e$.  Consequently two events
inside the same interval are indistinguishable from one event of
their combined size placed anywhere in that interval, and an event's
date becomes identified only when a listed expiry separates it from
its neighbours.
\end{proposition}

\begin{proof}
The clock \eqref{eq:clkclock} at a quoted expiry $t_i$ is
$365\,t_i+\sum_{t_e\le t_i}N_e$: it reads the calendar only through
cumulative sums up to quoted dates.  Two calendars with the same
per-interval totals give the same sums at every $t_i$, hence the same
$\tau(t_i)$, the same $\Delta\tau_i$, and the same value of every
quantity listed.  The date within an interval never enters.
\end{proof}

So the solver's placement of its candidate --- mid-interval, by
convention --- is exactly that: a convention, recorded in the
contract table, not an estimate.  When the report date is public,
type it in; the \emph{size} is what the board can be asked for.

What can the size estimator deliver?  \Cref{fig:clkident} is the
planted-event study: a flat clock
volatility, one event of known size, solve
\eqref{eq:clkobj}, compare.  Panel~(a) plots recovered against
planted size on three boards.  On a dense board --- expiries from a
week to a year, listed like a liquid single name's --- at $40\%$
volatility, the curve is the diagonal to the eye: recovery within
\MacClkIdentShrinkStrongD\ days at every planted size of two days
and up.  On the same board at index-like $20\%$ volatility, the
curve runs parallel to the diagonal but below it, short by up to
\MacClkIdentShrinkWeakD\ days.  And on a coarse quarterly board at
$20\%$, the curve is \emph{zero} until the planted event exceeds
\MacClkIdentBlindD\ days, then rises roughly parallel, several days
short.  (On every board the toe of the curve --- planted events
under a day or so --- reports as exactly zero; the last paragraph of
the derivation below says why.)  Three behaviours, one mechanism,
and it is worth deriving because it is the cost of the
sparsity charge.

Work on the study's configuration: one live candidate $N$ (the
others parked at zero, where the solve indeed leaves them)
against a planted truth $N^\ast$.  Two facts start the derivation.
On this board the flatness part of \eqref{eq:clkobj} is
\emph{minimized exactly at the truth}: the generator's ladder is flat
by construction, so installing $N=N^\ast$ restores a perfectly flat
$\mathcal{F}$ and no other value does better.  And near its minimum a
smooth function is a parabola: value flat at $N^\ast$, some curvature
$J''$, giving $\tfrac12 J''\,(N-N^\ast)^2$.  Adding the per-day
charge $(\lambda_{\rm sparse}/365)\,N$ tilts the parabola, and
setting the derivative to zero moves its minimum:
\[
J''(\widehat N-N^\ast)+\frac{\lambda_{\rm sparse}}{365}=0
\qquad\Longrightarrow\qquad
\widehat N=N^\ast-\frac{\lambda_{\rm sparse}}{365\,J''}.
\]
Every sparsity-charged estimator shrinks; the shortfall is the
charge over the curvature.  And the curvature is where the market
enters.  The candidate moves the objective only through its own
interval's forward variance, at rate
\[
\frac{\partial\mathcal{F}_i}{\partial N_i}
=-\frac{\Delta w_i}{\Delta\tau_i^{2}}\cdot\frac{1}{365}
=-\frac{\mathcal{F}_i}{365\,\Delta\tau_i},
\]
so the flatness curvature is of the order of the square of that
rate, $J''\sim\big(\mathcal{F}_i/(365\,\Delta\tau_i)\big)^2$, up to
the count of squared terms touched.  Since $\mathcal{F}_i$ is of
order $\sigma^2$, the shrinkage scales like
\[
N^\ast-\widehat N
\;\sim\;\lambda_{\rm sparse}\,
\frac{\Delta\tau_i^{2}}{\sigma^{4}}
\]
--- inversely with the \emph{fourth} power of the board's volatility,
and directly with the \emph{square} of the interval's width, with the
day-count conversions and term counts absorbed into the
proportionality (the study's measured pair,
\MacClkIdentShrinkStrongD\ against \MacClkIdentShrinkWeakD\ days,
carries rounding at the small end).  Halve the volatility and the
shrinkage grows roughly sixteenfold; stretch the bracketing interval
from a month to a quarter and it grows another order of magnitude.

One mechanical fact completes the picture.  The estimate is
constrained to $N\ge0$ and solved events below half a day are
dropped, so whenever the shifted optimum $\widehat N$ lands at or
below that floor, the solver reports \emph{exactly zero} --- the flat
toes at the left of panel~(a), and the walls of panel~(b), are this
clipping in action, not a failure to compute.  Panel~(b) shows the
consequence directly: fix the planted event at 2 days and sweep the
board's volatility.  The dense board starts seeing the event at
about \MacClkIdentWallWeeklyPct\% volatility; the quarterly board not
until about \MacClkIdentWallQuarterlyPct\%.  Below the wall the
flatness gain a small event offers cannot clear the per-day charge,
and the solver --- correctly, by its own lights --- says nothing.
One more line of the study completes the picture: handed a perfectly
flat ladder, the solver installs exactly \MacClkIdentFlatDays\ days.  A
sparsity prior's virtue is precisely that silence is its default.

\begin{figure}[htbp]
\centering
\paperfig{\textwidth}{fig_clk_ident.pdf}
\caption{The planted-event study.  (a)~Recovered vs planted size:
near-exact on the dense $40\%$ board (shrinkage
$\le\MacClkIdentShrinkStrongD$\,d for planted sizes of 2\,d and up),
shrunk by up to \MacClkIdentShrinkWeakD\,d at $20\%$, and blind below
\MacClkIdentBlindD\,d on a quarterly $20\%$ board --- the
$\lambda_{\rm sparse}\,\Delta\tau^2/\sigma^4$ law at work, with the
$N\ge0$ clipping flattening every curve's toe.  (b)~A fixed 2-day
event against board volatility: the dense board's materiality wall
sits near \MacClkIdentWallWeeklyPct\%, the quarterly board's near
\MacClkIdentWallQuarterlyPct\%.}
\label{fig:clkident}
\end{figure}

We now know the estimator's character.  It reads exactly the kinks
the clock was built to explain; it can only pull hot intervals down;
it shrinks what it reports and stays silent below a wall it computes
from the board's own volatility and spacing; and it cannot date an
event more finely than the listing brackets it.  What it has never
been told is \emph{why} an interval runs hot.  The last section runs
it on the real frozen boards and confronts that limit directly.
